\section{Will and Love}\label{sec:will}

From the angelic intellect Aquinas passes to the things that pertain to
the will of the angels (\latin{de his quae pertinent ad voluntatem
Angelorum}), considering first the will itself (q.~59) and then its
movement, which is love or dilection (\latin{de motu eius, qui est amor
sive dilectio}), for every act of the appetitive power is derived from
love or dilection (\latin{nam omnis actus appetitivae virtutis ex amore
seu dilectione derivatur}) (\ST{I q.~59 pr.}; \ST{I q.~60 pr.}).
The two questions carry forward the results of the treatise on knowledge
(\S\ref{sec:knowledge}): because the angel knows the universal aspect of
goodness by intellect, it loves; because its knowing is not discursive,
its choosing is not deliberative. The Church teaches that the angels, as
purely spiritual creatures, ``have intelligence and will: they are
personal and immortal creatures'' (CCC~330).

\subsection{The Will of the Angels (\ST{I q.~59})}\label{sec:q59}

The four articles ask whether there is will in the angels, whether the
will is their nature or their intellect, whether there is free choice in
them, and whether they have an irascible and a concupiscible appetite
(\ST{I q.~59 pr.}).

\paragraph{Will in the angels (a.~1).} Aquinas cites Augustine (\emph{De
Trinitate} X.11--12) for the \emph{sed contra}: the image of the Trinity
is found in the mind according to memory, understanding, and will, and
that image is in the angelic mind as well as the human. The corpus argues
from the grades of inclination. Every nature is inclined to good, but in
different ways: things without knowledge have \latin{appetitus
naturalis}; sense has \latin{appetitus sensitivus}, which reaches only
particular goods; the intellect, which apprehends goodness itself,
``inclined towards good in general. Such inclination is termed `will.'''
``Accordingly, since the angels by their intellect know the universal
aspect of goodness, it is manifest that there is a will in them.'' The
first objection had argued that the will is ``in the reason,'' and the
angels are above reason. Aquinas answers that reason and intellect differ
only in manner of knowing: reason reaches the universal by discourse,
intellect ``by simple intuition'' (\latin{simplici intuitu}); the object
proposed to the appetite is the same, so ``in the
angels, who are purely intellectual, there is no appetite higher than the
will'' (ad~1). Nor is will excluded because the angels are perfect and
unmoved: appetite is named from seeking what is not yet had, but it
reaches also to the good possessed (ad~2); and to will is a ``movement''
only in the sense in which the act of a perfect thing is so called
(ad~3).

\paragraph{Will distinct from nature and intellect (a.~2).} ``In the
angels the will is a special faculty or power, which is neither their
nature nor their intellect.'' The essence of a thing is contained wholly
within it, so what reaches beyond it to an extrinsic good must be
something superadded; only in God, in whom universal being and universal
good are contained, are intellect and will identical with essence --- in
no creature. Intellect and will differ by the direction of their
respective acts: knowledge is by the presence of the known within the
knower, while ``the will goes out to what is beyond it, according as by
a kind of inclination it tends, in a manner, to what is outside it.''
The \emph{sed contra} is simpler still: the angelic will regards good
things only, the intellect regards good and evil alike, ``for they know
both.'' The replies rest the distinction on the formal objects: the true
and the good are convertible in reality, but their diversity of aspect
``suffices for the difference of intellect from will'' (ad~2).

\paragraph{Free choice in the angels (a.~3).} The \emph{sed contra}
argues a fortiori from human dignity: free choice is part of man's
dignity, and the angels' dignity surpasses man's. The corpus defines the
matter. Some things act without judgment, as the arrow toward the
target; irrational animals act by a judgment ``not a free one, but
implanted by nature''; ``only an agent endowed with an intellect can act
with a judgment which is free,'' because it apprehends the common aspect
of goodness. ``Consequently, wherever there is intellect, there is
free-will'' (\latin{ubicumque est intellectus, est liberum arbitrium});
therefore there is \latin{liberum arbitrium} in the angels, ``in a
higher degree of perfection than in man.'' Damascene had taught the same
in the East: the angel's nature is ``rational, and intelligent, and
endowed with free-will, changeable in will,'' and precisely because
created it is ``changeable, having power either to abide or progress in
goodness, or to turn towards evil'' (\emph{De fide orthodoxa} II.3). The
first objection presses Aristotle's account of choice as desire
following counsel, and counsel as inquiry --- of which the undiscursive
angel is incapable. Aquinas answers that this describes human choice:
``there is choice in the angels, yet not with the inquisitive
deliberation of counsel, but by the sudden acceptance of truth''
(\latin{per subitam acceptionem veritatis}) (ad~1). The intellect's
determination to natural truths is no imperfection, so neither is the
will's determination toward what is above the angel (ad~2); and free
choice exists ``in a nobler manner'' in the higher angels, though liberty
as the removal of compulsion admits no degrees (ad~3).

\paragraph{No irascible or concupiscible appetite in the angels (a.~4).}
The first objection is Dionysius, who says of the demons that there is
in them ``unreasonable fury and wild concupiscence''; in Parker's
translation of the passage, ``what is evil in demons? An irrational
anger----a senseless desire----a headlong fancy'' (\emph{DN} 4.23). The
\emph{sed contra} is Aristotle: the irascible and concupiscible belong
to the sensitive part, which does not exist in angels. Aquinas
determines that ``the intellective appetite is not divided into
irascible and concupiscible; only the sensitive appetite is so
divided.'' Powers are distinguished by the formal aspect of their
objects, and the will's object is good under the common aspect of
goodness, which is not split into the arduous and the delectable; hence
the angelic appetite ``remains undivided; and it is called the will''
(\latin{remanet indivisus; et vocatur voluntas}). Fury and concupiscence
are said of the demons ``metaphorically \dots\ on account of the
resemblance in the effect'' (ad~1). Love and joy, as passions, belong to
the concupiscible; but as simple acts of the will they belong to the
intellective part, ``in this sense to love is to wish well to anyone''
(\latin{amare est velle bonum alicui}) --- and none of these is said of
the angels by way of passion, as Aquinas notes, citing Augustine
(\emph{De civitate Dei} IX) (ad~2). Charity and hope are in the will
itself, not in a sensory appetite; temperance and fortitude are said of
the angels only according as their will is conformed and firm in the
divine will (ad~3).

\angeldossier{\ST{I q.~59 aa.~1--4}; parallels \emph{Summa contra
gentiles} II.47--48.}
 {Church teaching: the angels have intelligence and will and are
 personal, immortal creatures (CCC~330). Common teaching: they have free
 choice (a.~3; Damascene, \emph{De fide orthodoxa} II.3). Thomist
 position: the will is a power distinct from nature and intellect (a.~2);
 angelic choice is without deliberative inquiry (a.~3); the intellective
 appetite is undivided (a.~4).}
 {Augustine, \emph{De Trinitate} X.11--12 (\emph{sed contra} of a.~1)
 and \emph{De civitate Dei} IX (a.~4 ad~2); Dionysius, \emph{DN} 4.23
 (obj.~1 of a.~4); Aristotle, \emph{De anima} III and \emph{Nicomachean
 Ethics} III; Damascene, \emph{De fide orthodoxa} II.3; CCC~330.}
 {The Dionysian attribution of fury and concupiscence to the demons
 (\emph{DN} 4.23), taken by Aquinas metaphorically (a.~4 obj.~1, ad~1);
 the rationalist placement of the will ``in the reason'' above which the
 angels would stand (a.~1 obj.~1), answered by the unity of the
 appetite's object in reason and intellect.}

\subsection{The Love or Dilection of the Angels (\ST{I q.~60})}\label{sec:q60}

Love is the first movement of the will and the root of every other; the
question asks whether there is natural love in the angels, and love of
choice, and how each bears on self, on the other angels, and on God
(\ST{I q.~60 pr.}).

\paragraph{Natural love (a.~1).} The \emph{sed contra} rests on
Augustine (\emph{De Trinitate} X.1--2): nothing is
loved unless first known, and the angels have natural knowledge,
therefore natural love. The corpus argues from the structure of nature
itself. ``It is common to every nature to have some inclination; and
this is its natural appetite or love.'' What belongs to nature as prior
is preserved in what follows, and nature is prior to intellect; so in
the intellectual nature there is a natural inclination ``coming from the
will,'' \latin{dilectio naturalis} seated in the angelic will itself.
The objection from Dionysius (\emph{DN} 4), that natural love is
contradistinguished from intellectual love, is answered: intellectual
love is divided only against the merely natural love of things without
sense or intellect (ad~1) --- Dionysius's own gradation runs from ``the
intellectual and rational,'' who aspire ``through knowledge,'' down to
things bereft of perception and things that merely are (\emph{DN} 4.4).
Natural love does not make the angel passive, since the Author of nature
implanted the inclination and the angel still acts by free will (ad~2);
and natural love is always right, for ``natural love is nothing else
than the inclination implanted in nature by its Author'' (\latin{amor
naturalis nihil aliud sit quam inclinatio naturae indita ab auctore
naturae}) --- though its rectitude is other than the rectitude of
charity, and is perfected by it (ad~3).

\paragraph{Elective love (a.~2).} The \emph{sed contra} is drawn from
merit: we neither merit nor demerit by purely natural acts, yet by their
love the angels merit or demerit (\latin{naturalibus neque meremur
neque demeremur}). ``There exists in the angels a natural love, and a
love of choice'' (\latin{dilectio electiva}), and the natural love is
the principle of the elective, as nature is principle to what follows
it. The analogy is exact: as the intellect knows principles naturally
and concludes from them, so the will tends naturally to the last end ---
``every man naturally wills happiness'' --- and chooses the means for
the end's sake. The angel's intellect, being perfect, has no discursive
knowledge; but its will has both loves, because it is no imperfection to
seek one thing naturally as end and another by choice as ordained to
that end. Aquinas closes by deferring all that is above nature ---
``nature is not the sufficient principle thereof'' --- to the treatment
of grace and glory (a.~2; \S\ref{sec:q62}). The first objection, again
from \emph{DN} 4, had set rational love against intellectual: not every
elective love follows discursive reasoning, but only human choice
(ad~1).

\paragraph{Love of self (a.~3).} The \emph{sed contra} is Aristotle:
``Love for others comes of love for oneself'' (\emph{Nic.\ Eth.} IX.8).
Good is found both in subsisting things and in accidents, so a thing is
loved in two ways: as a subsisting good, to which we wish well --- the
love called ``friendship'' (\latin{amor amicitiae}) --- and as a good to
be had, the love called ``concupiscence'' (\latin{amor
concupiscentiae}). Everything naturally seeks its own good and
perfection, and this is to love self; hence ``angel and man naturally
love self,'' and each also loves self with the love of choice, ``in so
far as from choice he wishes for something which will benefit himself.''
The two loves do not divide the object but the aspect (ad~1). Against
Dionysius's description of love as ``a uniting and a binding power'' ---
in Parker's rendering, love is ``of a power unifying, and binding
together, and mingling pre-eminently in the Beautiful and Good''
(\emph{DN} 4.12) --- Aquinas answers that ``to be one is better than to
be united,'' so there is more oneness in the love directed to self;
Dionysius used those terms ``to show the derivation of love from self
to things outside self'' (ad~2). Love is a movement that abides within
the lover and can be reflected back upon him, as knowledge is reflected
upon the knower (ad~3).

\paragraph{Love of the other angels (a.~4).} The \emph{sed contra} is
Sir 13:19: ``Every beast loveth its like'' (Douay). The corpus builds on
natural unity. ``What is one with a thing, is that thing itself'': what
is one with a thing by natural union it loves with natural love, as a
man loves his kinsman; what is one with it in species it loves ``in so
far as it loves its own species.'' ``So then, it must be said that one
angel loves another with natural affection, in so far as he is one with
him in nature.'' The ``as himself'' marks likeness, not equality: on the
side of the knower, an angel knows himself by his essence and another
not so, and loves himself by his own will, not by the other's will
(ad~1); and since natural love rests on natural unity, the angel loves
more what is numerically one with himself, yet ``it is natural for him
to have a like love for another as for himself, in this respect, that as
he loves self in wishing well to self, so he loves another in wishing
well to him'' (ad~2). Even the fallen angels keep this natural love: it
cannot be stripped from them ``without their still retaining a natural
affection towards the good angels, in so far as they share the same
nature with them. But they hate them, in so far as they are unlike them
according to righteousness and unrighteousness'' (ad~3).

\paragraph{Love of God above self (a.~5).} Aquinas first reports a
contrary opinion: some held that the angel loves God more than self only
in certain respects --- seeking the divine good for himself, and
wishing that God be God --- ``but absolutely speaking, out of the
natural love he loves himself more than he does God.'' The falsity of
this appears from the natural order: everything that naturally belongs
to another ``is principally, and more strongly inclined to that other to
which it belongs, than towards itself,'' as the hand exposes itself for
the whole body without deliberation; and reason copies nature, for the
virtuous citizen faces death for the commonwealth. ``Consequently, since
God is the universal good, and under this good both man and angel and
all creatures are comprised, because every creature in regard to its
entire being naturally belongs to God, it follows that from natural love
angel and man alike love God before themselves and with a greater love''
(\latin{naturali dilectione etiam Angelus et homo plus et principalius
diligat Deum quam seipsum}). Otherwise natural love would be perverse,
and ``would not be perfected but destroyed by charity.'' The replies
refine the principle: where one thing is the whole cause of another's
existence and goodness, that one is naturally loved more than self
(ad~1); God is loved for his own sake, though the capacity to love him
comes of depending on him (ad~2). God as the universal good ``is loved
by everything with natural'' love; ``so far as He is the good which of
its very nature beatifies all with supernatural beatitude, He is loved
with the love of charity'' (ad~4). And because the divine essence and
the universal good are one, whoever sees that essence cannot but love
him; those who do not see it know him by particular effects sometimes
contrary to their will, and so are said to hate him (ad~5). Augustine's
account of the two cities is read through this lens: the angels who
stood ``steadfastly continued in that which was the common good of all,
namely, in God Himself,'' while the others, ``enamored rather of their
own power, as if they could be their own good, lapsed to this private
good of their own'' (\emph{De civitate Dei} XII.1).

\angeldossier{\ST{I q.~60 aa.~1--5}.}
 {Common teaching: the angels love God above self with natural love
 (a.~5) and merit or demerit by a love of choice (a.~2); natural love is
 of itself upright (a.~1); each angel loves himself and his like
 (aa.~3--4).}
 {Sir 13:19 (\emph{sed contra} of a.~4); Augustine, \emph{De Trinitate}
 X.1--2 (\emph{sed contra} of a.~1) and \emph{De civitate Dei} XII.1, 9
 and XIV (objections of a.~5); Dionysius, \emph{DN} 4.4, 4.12
 (objections of aa.~1--3); Aristotle, \emph{Nic.\ Eth.} IX.8 (\emph{sed
 contra} of a.~3).}
 {The unnamed opinion (\latin{quidam}) that the angel loves God more
 than himself only secundum quid, and himself more absolutely, which
 Aquinas rejects as making natural love perverse (a.~5); Dionysius's
 ``uniting and binding'' account of love pressed against self-love
 (a.~3 obj.~2).}
